% !TEX root = userguide.tex

\chapter{Connections}
\pgst{tfemch}

\label{ch:connections}

\def\chaptertitle{Connections}
\def\headdate{18th February 2015}

In this chapter it is shown how to handle connections in \tfem.

\section{Use of connections}

Connections are used to add `direct stiffness' between degrees of freedom in
one part of the mesh and another part of the mesh. By `direct stiffness' we
mean that no additional unknowns, such Lagrange multipliers in constraints
(see Chap.~\ref{ch:constraints}) or from additional meshes, are used. Only
the sparsity of the existing matrix changes due to adding of stiffness between
unknowns that would otherwise have zero direct interaction.
Examples:
\begin{enumerate}
\item Two-phase flow with interface tension. In Fig.~\ref{fig:interface} a
fluid mesh is shown with an additional mesh for the interface between two
fluids. The unknowns are the velocity/pressure in the fluid nodes and the
position vector in the interface nodes.
\begin{figure}[ht!]
   \begin{center}
      \includegraphics[scale=0.7]{interface}
   \end{center}
   \caption{Mesh consisting of fluid elements and interface elements}
    \label{fig:interface}
\end{figure}
The effect of the interface tension is a distributed force contribution to the
momentum balance of the fluid. This force depends on the position of the
interface. The position of the interface depends on the local velocity at the
interface. In the standard finite-element matrix there is no coupling between
the velocity and the position vectors at the interface. By defining a
connection along the interface, the coupling can be achieved at the matrix
level. This opens up the possibility to compute the movement of the droplet
using fully implicit schemes avoiding very small time steps.

\item Coupling of domains of different resolution using the discontinuous
Galerkin method at the interface between the domains.

\item Coupling of two fluid domains with a viscous slip law (or other interface
rheology) in between the fluids.

\end{enumerate}

For a connection there are always two sides. The first is the defining side
and the second is the side that is connected to the first.
The possibilities implemented in \tfem\ are
given in the following table:
\begin{center}
\begin{tabular}{llcc}
side1 & side2 & weak & collocation \\\hline
point1 & point2 & $-$ & $+$ \\
nodeset1 & nodeset2 & $-$ & $+$ \\
curve1 & curve2 & $+$ & $+$  \\
      & elementset2 & $+$ & $+$ \\
      & nodeset2 & $-$ & $+$ \\
surface1 & surface2 & $+$ & $+$  \\
        & elementset2 & $+$ & $+$ \\
      & nodeset2 & $-$ & $+$ \\
volume1 & volume2 & $+$ & $+$  \\
        & elementset2 & $+$ & $+$ \\
      & nodeset2 & $-$ & $+$ \\
elementset1 & elementset2 & $+$ & $+$  \\
object & second side of object & $+$ & $+$  \\
object & object2 & $+$ & $+$  \\
       & curve2 & $+$ & $+$  \\
       & surface2 & $+$ & $+$  \\
       & volume2 & $+$ & $+$  \\
       & elementset2 & $+$ & $+$ \\
       & nodeset2 & $-$ & $+$ 
\end{tabular}
\end{center}
Remarks:
\begin{itemize}
\item   
A connection
is distributed, but is always \emph{local}, i.e.\ along the
defining side (side1) there will be a connection between elements or nodes that
are connected locally to the elements or nodes on the connected side (side2).
\item A \emph{weak} connection connects the two sides element by element,
whereas a \emph{collocated} connection connects the two sides node for node.
\item There is a fundamental difference between using objects and the other
geometrical entities.
Objects do not define degrees of freedom. The degrees of freedom are
defined by the elements that are intersected by the object. The connection
involves the degrees of freedom in the
\emph{intersected elements} both for weak and collocated
connections. The other geometrical entities
do define degrees of freedom and the connection involves the degrees of
freedom in the \emph{elements} for weak connections and \emph{nodes} for
collocation connections.
\item For points and nodesets only collocation makes sense. Degrees of freedom
in nodes are connected. The sequence of the nodes in both nodesets is
important, since the connection is node for node.
\item Curves and surfaces represent element edges and faces. Connections
between curves or surfaces create connections between degrees of freedom on
element edges/faces (weak) or nodes (collocation).
\item Curves, surfaces and volume can also be connected to elementsets and
nodesets. The latter obviously only with collocation.
Note, that elementsets represent full mesh elements and not just
element edges or faces. 
\item Objects can be connected to itself (double-sided objects) or to another
object.
\item Objects can also be connected to curves, surfaces, volumes,
 elementsets and nodesets. The latter obviously only with collocation. 
\item Restrictions:
\begin{itemize}
   \item For a weak connection the number of elements must be the same on both
         sides.
   \item In addition, objects on both sides of a weak connection need to have
         the same element type and the same number of integration points.
   \item For a collocated connection the number of nodes must be the same on
         both sides.
\end{itemize}
Note, that it is allowed to use different element types on both sides of the
connection, except for objects on both sides of the connection.
\end{itemize}

\section{Implementation}

\label{sec:connimpl}

For geometrical entities that are not objects (on side1), the system
matrix is built element wise (weak) or
nodal point wise (collocation). 
Note, that for the weak case, the element matrix of the connection
involves degrees of freedom of two elements, one on each
side of the connection. For the collocation case,
the element matrix of the connection
involves degrees of freedom of two nodes, one on each 
side of the connection\footnote{Be careful here, since
it is tricky to define the correct stiffness at the nodes. The
stiffness is more like discrete `springs' attached to the nodes.}.

For objects (on side1), the system matrix is built integration
point wise (weak) or
nodal point wise (collocation).
If the second side is not an object, but another
geometrical entity (curve,\ldots), keep in mind that
the system matrix is still built
integration point wise (on the object of side1) for the weak case. 
Note, that for objects
the element matrix of the connection
involves degrees of freedom of intersected elements 
for both the weak and collocated case.

At the call of the element subroutine, the user must supply
four element matrices ($\mat S_{11}$, $\mat S_{12}$, $\mat S_{21}$, 
$\mat S_{22}$) and two element vectors ($\col f_1$, $\col f_2$) that are
added to the system. The full matrix and vector that is added is
\[
 \begin{pmatrix}
    \mat S_{11} & \mat S_{12} \\
    \mat S_{21} & \mat S_{22}
 \end{pmatrix},
\quad
 \begin{pmatrix}
    \col f_{1}  \\
    \col f_{2} 
 \end{pmatrix}
\]
where the indices 1 and 2 denote the unknowns 
from side 1 and side 2, respectively. For objects, the weight of the
integration rule at the current integration point
needs to be taken into account on element level.

The heading of the element subroutine defining a connection is given by
\begin{listing}{1}
   
  subroutine elemsub ( mesh, problem, conn, elem, node, matrix, vector, &
    first, last, coefficients, oldvectors, elemmat11, elemmat12, &
    elemmat21, elemmat22, elemvec1, elemvec2 )
    type(mesh_t), intent(in) :: mesh
    type(problem_t), intent(in) :: problem
    integer, intent(in) :: conn, elem, node
    logical, intent(in) :: matrix, vector, first, last
    type(coefficients_t), intent(in) :: coefficients
    type(oldvectors_t), intent(in) :: oldvectors
    real(dp), intent(out), dimension(:,:) :: elemmat11, elemmat12, &
      elemmat21, elemmat22
    real(dp), intent(out), dimension(:) :: elemvec1, elemvec2

     ....

  end subroutine elemsub

\end{listing}
with\\[2ex]
\begin{tabular}{ll}
\texttt{mesh, problem} & mesh and problem structure as in the heading of
                           \texttt{build\_system\_connection}\\
\texttt{conn} & connection number \\
\texttt{elem} & element number in the geometry or object
                for a weak connection \\
\texttt{node} & node or integration point number within this connection\\
\texttt{matrix, vector} & logicals defined by \texttt{buildmatrix} and
                         \texttt{buildvector} in the heading \\
                        & of \texttt{build\_system\_connection}.\\
\texttt{first, last} & logicals indicating whether the current call
    is the first/last \\ & in this connection. \\
\texttt{coefficients}, & coefficient and oldvectors as supplied
        in the heading of \\
\texttt{oldvectors} & \hfill \texttt{build\_system\_connection}.\\
\texttt{elemmat11} & matrix $\mat S_{11}$. \\
\texttt{elemmat12} & matrix $\mat S_{12}$. \\
\texttt{elemmat21} & matrix $\mat S_{21}$. \\
\texttt{elemmat22} & matrix $\mat S_{22}$. \\
\texttt{elemvec1} & vector $\col f_{1}$. \\
\texttt{elemvec2} & vector $\col f_{2}$.
\end{tabular}\\[2ex]

In addition to \texttt{elemsub}, two alternative element interfaces,
\texttt{elemsub1} and \texttt{elemsub2}, that have an extended heading with
\texttt{elgrp1}, \texttt{elem1}, \texttt{elgrp2}, \texttt{elem2} and
\texttt{nodenr1}, \texttt{nodenr2}, respectively.

Multiple connections can be defined, However, 
you need to be careful:
\begin{enumerate}[a.]
\item   If connections have identical values for BOTH \texttt{physq1} and
        \texttt{physq2} (\texttt{physq1} may be different from \texttt{physq2})
        or have different values for BOTH \texttt{physq1} and \texttt{physq2},
        there are no complications. In the first case the connections are simply
        merged into basically one single connection. In the second case
        the degrees of freedom are separated by default.
\item   If NOT a. then be sure that the connections are FULLY SEPARATED, i.e.
        the element nodes involved on both sides of the connection MUST NOT
        OVERLAP!
\end{enumerate}


\section{Examples of problems using connections}

The use of connections is very similar to the use of constraints:
\begin{enumerate}
   \item Connections are defined using the routine \texttt{define\_connection}
     before the problem definition.
   \item The structure of the matrix is added using the routine
     \texttt{create\_sysmatrix\_structure\_connection} before the final call
    to \texttt{finalize\_sysmatrix\_structure}.
   \item The system matrix contribution of the connections is build using
     \texttt{build\_system\_connection}.
\end{enumerate}

In Secs.~\ref{sec:stokes20} and \ref{sec:stokes40} connections are
 used to add a viscous
slip layer at an interface between two fluids in a cavity flow.
In the next section
(Sec.~\ref{sec:stokes41}) the one-sided connection is explained.
In Sec.\ref{sec:stick_slip}, example \texttt{stick\_slip5.f90}, uses a one-sided
connection to create a developed entry flow profile.

\subsection{A Stokes problem on a unit square with a one-sided connected
            boundary condition}
\label{sec:stokes41}

This example is similar to \texttt{stokes2.f90} (see
Sec.~\ref{sec:stokes2}), but the constraint for the periodical boundary
conditions has been replaced with a collocated connection between the out- en inflow.
The source is in the file \texttt{stokes41.f90}
and can be obtained by typing at the
command line:
\begin{quote}
   \texttt{getexample stokes}
\end{quote}
This is a rather artificial problem, where the outflow is now a free surface 
and the velocities at the outflow are fed back to the inflow. Note, that this is different
from periodical boundary conditions. There is no stress continuity between in and out,
only continuity of velocity. The flow is generated with a constraint for an
imposed flow rate at the outflow boundary.
A velocity vector plot is shown in Figure~\ref{fig:vel41}.
\begin{figure}[ht!]
   \begin{center}
   \includegraphics[scale=0.3]{velocity_stokes41}
   \end{center}
   \caption{Velocity vector plot for the one-sided connection problem. The
inflow velocity profile is set equal to the outflow velocity profile.}
    \label{fig:vel41}
\end{figure}

The problem is solved basically as follows. First the system is built,
including the constraint for the imposed flow rate, but  
without the connection.
Then the rows in the system matrix and the entries in the right-hand side
belonging to the degrees of freedom 
on the inflow curve
are set to zero. These rows are then filled by 1's and -1's using the connection to
exactly represent the following equation node for node:
\[
   \vek u_\text{in}-\vek u_\text{out}=\vek 0
\]
Note, that the system matrix is non-symmetric.

Remark:
\begin{itemize}
   \item Constraints should \emph{not} involve the degrees of freedom of which
the rows in the matrix and right-hand side have been set to zero.
\end{itemize}

